\chapter{Criterio de protocolización mínima}

\section{Protocolización disciplinada}
\begin{definition}
Se llama protocolización mínima a la explicitación ordenada de pasos, condiciones y criterios de uso mediante los cuales una configuración, familia o traducción puede ser tratada de forma repetible dentro del canon intermedio.
\end{definition}

\begin{observation}
La protocolización mínima no equivale todavía a una metodología completa ni a una auditoría final. Su función es fijar un primer nivel de repetibilidad disciplinada.
\end{observation}

\section{Condición mínima de protocolo}
\begin{definition}
Existe protocolización mínima cuando pueden declararse de manera legible:
\begin{enumerate}
  \item el objeto tratado;
  \item su dependencia canónica;
  \item el nivel de tratamiento en que opera;
  \item las condiciones mínimas de entrada;
  \item las condiciones mínimas de salida o fallo.
\end{enumerate}
\end{definition}

\begin{observation}
No se exige todavía una normalización exhaustiva de todos los pasos posibles. Basta con una declaración suficiente para evitar tratamiento arbitrario o irrepetible.
\end{observation}

\section{Protocolización y configuración}
\begin{definition}
Toda configuración operativa puede protocolizarse mínimamente cuando su función dentro del canon puede expresarse mediante reglas de uso no contradictorias.
\end{definition}

\begin{observation}
Esta cláusula evita que entren al tramo protocolizable piezas cuya función siga siendo puramente verbal o indeterminada.
\end{observation}

\section{Protocolización y familia}
\begin{definition}
Una familia operativa admite protocolización mínima cuando sus miembros comparten un núcleo de tratamiento suficientemente común como para sostener reglas de uso compatibles.
\end{definition}

\begin{observation}
La protocolización de una familia no exige uniformidad absoluta entre todos sus miembros. Exige solo una base común suficiente para tratamiento disciplinado.
\end{observation}

\section{Protocolización y traducción}
\begin{definition}
Una traducción entre niveles es protocolizable cuando puede declararse qué se conserva, qué se pierde y bajo qué condiciones la pieza traducida sigue siendo tratable en el nivel de llegada.
\end{definition}

\begin{observation}
Esta cláusula impide que la traducción se use como desplazamiento opaco entre niveles.
\end{observation}

\section{Condición mínima de fallo}
\begin{definition}
Se considera fallida una protocolización mínima cuando no puede declarar de forma legible objeto, dependencia, nivel, condiciones de entrada o condiciones de salida sin contradicción inmediata.
\end{definition}

\begin{observation}
Esta condición protege al canon intermedio frente a protocolos aparentes que no alcanzan repetibilidad real.
\end{observation}

\section{Cláusula de alcance}
\begin{observation}
El criterio de protocolización mínima fija una primera gramática de repetibilidad disciplinada. No constituye todavía una teoría completa de auditoría, validación o trazabilidad final.
\end{observation}

